\chapter{Kết luận và hướng phát triển}
\label{ch:ketluan}

\begin{muctieu}
Chương này tổng kết kết quả, bài học rút ra và các hướng phát triển khả thi
từ những hạn chế đã phân tích.
\end{muctieu}

\section{Kết luận}

Đồ án đã xây dựng một hệ thống mobile hoàn chỉnh cho khám phá công thức và
nhận diện món Việt. Ứng dụng Flutter kết nối FastAPI, database SQLite và model
ViT-L/16; nhãn model được chuẩn hóa rồi ánh xạ tới 38 món có công thức, ảnh,
nguồn và giấy phép. Ngoài luồng AI, hệ thống hỗ trợ tìm kiếm, lọc, tài khoản,
yêu thích, lịch sử và phản hồi.

Đóng góp quan trọng của đồ án là kết nối ba loại dữ liệu vốn tách rời: ảnh
nhận diện, nhãn model và catalog công thức. Quy trình audit giúp tránh sử
dụng ảnh train làm ảnh giao diện, tránh công thức không rõ nguồn và tránh
tuyên bố hiệu năng model dựa trên một mẫu quá nhỏ. Đánh giá đầy đủ được thực
hiện trên 5.014 ảnh Test sạch thuộc 30 lớp có thể đối chiếu.

\section{Bài học rút ra}

Chất lượng dữ liệu ảnh ảnh hưởng trực tiếp tới độ tin cậy của đánh giá. Nếu
không loại ảnh trùng qua split, kết quả có thể lạc quan giả tạo. Với dữ liệu
công thức, có nội dung chưa đủ; mỗi ảnh và công thức còn cần metadata nguồn,
giấy phép và dấu vết thay đổi. Vì vậy phần audit và importer chiếm vai trò
không nhỏ hơn việc tích hợp model.

\begin{samepage}
Một model mạnh không tự tạo thành sản phẩm tốt. Hệ thống phải xử lý ảnh hỏng,
model chưa sẵn sàng, mất mạng, token hết hạn, nhãn không ánh xạ được và recipe
không hợp lệ. Tách UI, service, state, router và repository giúp từng trường
hợp có thể kiểm thử độc lập. Nếu thực hiện lại, việc xây dựng tiêu chí dữ liệu
và bộ test độc lập sẽ được ưu tiên ngay từ đầu.
\end{samepage}

Trí tuệ nhân tạo trong hệ thống đóng vai trò gợi ý, không phải kết luận tuyệt
đối. Confidence phải được trình bày minh bạch và người dùng cần có cơ chế gửi
phản hồi. Tương tự, công cụ AI hỗ trợ phát triển phải được kiểm chứng bằng mã
nguồn, test và nguồn tham khảo trước khi đưa vào báo cáo.

\section{Hướng phát triển}

\begin{itemize}
  \item \textbf{Đủ dữ liệu cho 38 nhãn:} xây dựng tập test có giấy phép và
        quy trình gán nhãn cho tám món còn thiếu, từ đó đánh giá đúng toàn bộ
        phạm vi production.
  \item \textbf{Fine-tune có kiểm soát:} huấn luyện candidate trên manifest
        sạch, so sánh trên test độc lập và chỉ thay production model khi có
        cải thiện được chứng minh.
  \item \textbf{Tối ưu suy luận:} thử nghiệm model nhỏ hơn, quantization,
        ONNX/TFLite hoặc GPU server để giảm thời gian và dung lượng.
  \item \textbf{Hạ tầng production:} chuyển sang PostgreSQL, rate limiter
        dùng kho chia sẻ, reverse proxy, HTTPS, log tập trung và monitoring.
  \item \textbf{Vòng phản hồi:} tổng hợp nhãn người dùng xác nhận thành hàng
        đợi kiểm duyệt, không đưa thẳng dữ liệu chưa kiểm chứng vào train.
  \item \textbf{Mở rộng nội dung:} bổ sung dinh dưỡng, dị ứng và biến thể
        vùng miền khi có nguồn chuyên môn phù hợp.
\end{itemize}
